\section{Related Work}

In this section, we will briefly review several lines of works closely related to ours, including general recommendation, sequential recommendation, and attention mechanism.


\subsection{General Recommendation}

Early works on recommendation systems typically use Collaborative  Filtering (CF) to model users' preferences based on their  interaction histories~\cite{Sarwar:www2001:ICF,Koren:ACF:2011}.
Among various CF methods, Matrix Factorization (MF) is the most popular one, which projects users and items into a shared vector space and estimate a user's preference on an item by the inner product between their vectors~\cite{Salakhutdinov:nips2007:PMF,Koren:2009:MFT,Koren:ACF:2011}.
Another line of work is item-based neighborhood methods \cite{Sarwar:www2001:ICF,Linden:ieee2003:ARI,Koren:kdd2008:FMN,Kabbur:kdd2013:FFI}.
They estimate a user's preference on an item via measuring its similarities with the items in her/his interaction history using a precomputed item-to-item similarity matrix.




Recently, deep learning has been revolutionizing the recommendation systems dramatically.
The early pioneer work is a two-layer Restricted Boltzmann Machines (RBM) for collaborative filtering, proposed by \citeauthor{Salakhutdinov:icml2007:RBM}~\cite{Salakhutdinov:icml2007:RBM} in Netflix Prize\footnote{\url{https://www.netflixprize.com}}. %


 One line of deep learning based methods seeks to improve the recommendation performance by integrating the distributed item representations learned from auxiliary information, \textit{e.g.}, text~\cite{Wang:kdd2015:CDL,Kim:recsys2016:CMF}, images~\cite{Wang:www2017:YIR,Kang:icdm2017:Visually}, and acoustic features~\cite{Oord:nips2013:Deep} into CF models. 
Another line of work seeks to take the place of conventional matrix factorization. 
For example, Neural Collaborative Filtering (NCF)~\cite{He:www2017:NCF} estimates user preferences via Multi-Layer Perceptions (MLP) instead of inner product, while AutoRec~\cite{Sedhain:www2015:AAM} and CDAE~\cite{Wu:wsdm2016:CDA} predict users' ratings using Auto-encoder framework.


\subsection{Sequential Recommendation}

Unfortunately, none of the above methods is for sequential recommendation since they all ignore the order in users' behaviors.

Early works on sequential recommendation usually capture sequential patterns from user historical interactions using Markov chains (MCs).
For example, \citeauthor{Shani:jmlr2005:MRS}~\cite{Shani:jmlr2005:MRS} formalized recommendation generation as a sequential optimization problem and employ Markov Decision Processes (MDPs) to address it. 
Later, \citeauthor{Rendle:WWW2010:FPM}~\cite{Rendle:WWW2010:FPM} combine the power of MCs and MF to model both sequential behaviors and general interests by Factorizing Personalized Markov Chains (FPMC).
Besides the first-order MCs, high-order MCs are also adopted to  consider more previous items~\cite{He:icdm16:Fusing,He:recsys2017:TR}. 
 
Recently, RNN and its variants, Gated Recurrent Unit (GRU)~\cite{Cho:emnlp2014:gru} and Long Short-Term Memory (LSTM)~\cite{Hochreiter:LSTM:NC1997}, are becoming more and more popular for modeling user behavior sequences~\cite{Hidasi:ICLR2016:gru4rec,Yu:sigir2016:DRM,Wu:wsdm2017:RRN,Donkers:recsys2017:SUR,Quadrana:recsys2017:PSR,Li:cikm2017:NAS,Hidasi:CIKM2018:RNN}.
The basic idea of these methods is to encode user's previous records into a vector (\textit{i.e.}, representation of user's preference which is used to make predictions) with various recurrent architectures and loss functions, including session-based GRU with ranking loss (GRU4Rec)~\cite{Hidasi:ICLR2016:gru4rec}, Dynamic REcurrent bAsket Model (DREAM)~\cite{Yu:sigir2016:DRM}, user-based GRU~\cite{Donkers:recsys2017:SUR}, attention-based GRU (NARM)~\cite{Li:cikm2017:NAS}, and improved GRU4Rec with new loss function (\textit{i.e.}, \texttt{BPR-max} and \texttt{TOP1-max}) and an improved sampling strategy~\cite{Hidasi:CIKM2018:RNN}.



Other than recurrent neural networks, various deep learning models are also introduced for sequential recommendation~\cite{Chen:wsdm2018:SRU,Tang:WSDM2018:caser,Liu:kdd2018:STAMP,Kang:ICDM2018:SAS}. %
For example, \citeauthor{Tang:WSDM2018:caser}~\cite{Tang:WSDM2018:caser} propose a Convolutional Sequence  Model (Caser) to learn sequential patterns using both horizontal and vertical convolutional filters.
\citeauthor{Chen:wsdm2018:SRU}~\cite{Chen:wsdm2018:SRU} and \citeauthor{Huang:sigir2018:ISR}~\cite{Huang:sigir2018:ISR} employ Memory Network to improve sequential recommendation.
STAMP captures both users' general interests and current interests using an MLP network with attention~\cite{Liu:kdd2018:STAMP}.

 

\subsection{Attention Mechanism}

Attention mechanism has shown promising potential in modeling sequential data, \textit{e.g.}, machine translation~\cite{Bahdanau:iclr2015,Vaswani:nips2017:Attention} and text classification~\cite{Yang:naacl2016:HAN}. 
Recently, some works try to employ the attention mechanism to improve recommendation performances and interpretability~\cite{Li:cikm2017:NAS,Liu:kdd2018:STAMP}.
For example, \citeauthor{Li:cikm2017:NAS}~\cite{Li:cikm2017:NAS} incorporate an attention mechanism into GRU to capture both the user's sequential behavior and main purpose in session-based recommendation.


\input{model_fig}

The works mentioned above basically treat attention mechanism as an additional component to the original models. %
In contrast, Transformer~\cite{Vaswani:nips2017:Attention} and BERT~\cite{Devlin:2018:BERT} are built solely on multi-head self-attention and achieve state-of-the-art results on text sequence modeling.
Recently, there is a rising enthusiasm for applying purely attention-based neural networks to model sequential data for their effectiveness and efficiency~\cite{Lin:iclr2017,Shaw:naacl2018:self,Liu:iclr2018:Generating,Radford:2018:GPT,Radford:2019:Language}.
For sequential recommendation, \citeauthor{Kang:ICDM2018:SAS}~\cite{Kang:ICDM2018:SAS} introduce a two-layer \textit{Transformer decoder} (\textit{i.e.}, Transformer language model) called SASRec to capture user's sequential behaviors and achieve state-of-the-art results on several public datasets.
SASRec is closely related to our work.
However, it is still a unidirectional model using a casual attention mask.
While we use a bidirectional model to encode users' behavior sequences with the help of Cloze task.

